\begin{proposition}\label{prop:all-octic-33}
Every rational effective divisor of degree $8$ on $X_1(33)$
supported away from the cusps is one of the five closed octic
divisors in the diamond orbit of
\[
 x^4-3x^3+3x^2-3x+1=0,\qquad
 y^2+(x^3-3x^2+x-3)y+2x^2-x+2=0.
\]
Consequently, if an elliptic curve $E$ over an octic number field $K$
has a point of order $33$, then $E(K)_{\mathrm{tors}}\simeq C_{33}$.
In particular, $C_{66}$, $C_{99}$ and $C_3\times C_{33}$ do not
belong to $\Phi(8)$.
\end{proposition}
\begin{proof}
Put $X=X_1(33)$, and choose a rational cusp $c_0$.
By \cite[Theorem~3.1(2), Theorem~4.13 and Table~2]{DEHMZ},
the full group $J_1(33)(\mathbb Q)$ is rational cuspidal with
invariants $[5,42373650]$.
An exact class-group computation over $\mathbb F_2$ gives
\[
 J_1(33)(\mathbb F_2)\simeq C_{15}\times C_{42373650}.
\]
The reductions of the rational cuspidal generators generate a subgroup
$H$ with invariants $[5,42373650]$. This is the full order of
$J_1(33)(\mathbb Q)$, so reduction on $J_1(33)(\mathbb Q)$ is
injective, including on its $2$-primary subgroup. The index of $H$
in the finite Jacobian is $3$.

The rational gonality of $X$ is $10$ by \cite[Table~1]{DvH}.
Thus distinct rational effective divisors of degree $8$ have distinct
classes in $\operatorname{Pic}^8(X)$: a nontrivial linear equivalence
would give a rational function of degree at most $8$.
It follows that two known rational divisors whose reductions have the
same class must be equal.

The rational cusp-orbit degrees are $1^{10},2^5,10^2$, while the
closed cusp degrees over $\mathbb F_2$ are $1^{10},2^5,5^2,10$.
We identify the $17$ rational cusp orbits by the valuations of the
integral Tate modular units $b$ and $b^4/\Delta$.
These units are invertible on the fine scheme $Y_1(33)$ over
$\mathbb Z[1/33]$. Their divisors have no vertical components and
are supported on the finite etale cusp boundary. Their cusp valuations
therefore commute with reduction at $2$.
The computed $17$ distinct signatures separate the rational orbits.

There are exactly $62055$ rational effective cuspidal divisors of
degree $8$. Their reductions give $62055$ distinct classes.
The complete finite-field enumeration, independently checked using
the modular-symbol Frobenius polynomial, gives noncuspidal place counts
\[
 (0,0,0,0,2,20,30,25)
\]
in degrees one through eight.
Every degree-eight reduction is therefore entirely cuspidal or
has a single noncuspidal place of degree $5$, $6$, $7$, or $8$,
with a cuspidal complement. All multiplicities in the cuspidal part
are included. The complete class tests are
\[
\begin{array}{c|r|r|r}
\text{noncuspidal part}&\text{divisors tested}&
 \text{class in }H&\text{rational effective cuspidal classes}\\\hline
\text{empty}&62595&62055&62055\\
(5)&540&0&0\\
(6)&1200&0&0\\
(7)&300&100&0\\
(8)&25&5&0
\end{array}
\]
The numbers of cuspidal complements of degrees $0,1,2,3$ are
$1,10,60,270$, respectively. No mixed noncuspidal support is possible,
since its degree would be at least $10$.

The $100$ classes in the degree-seven row are reductions of known
rational divisors of degree $8$: add any of the ten rational cusps
to each of the ten diamond translates of the septic point
\[
 a^7-6a^6+9a^5-5a^3-a+1=0,\qquad
 y=\frac{-3a^6+16a^5-17a^4-9a^3+11a^2+3a+4}{3}.
\]
We reconstruct this point in characteristic zero and verify its order
$33$ afresh. Only the existence of these reducible degree-eight
divisors is used.
The remaining five classes are precisely the reductions of the five
diamond translates of the octic point in the statement.
Its absolute field polynomial is
\[
 t^8-15t^7+95t^6-331t^5+741t^4-1100t^3+1009t^2-688t+359,
 \qquad t=x+y.
\]
We check its irreducibility, reconstruct the original coordinates,
and verify exact order $33$ over this octic field.

Thus every possible reduction class is either a rational effective
cuspidal class, one of the $100$ septic-plus-cusp classes, or one
of the five known octic classes. Injectivity of reduction identifies
the corresponding global classes, and gonality $10$ identifies the
effective divisors. A divisor supported away from the cusps cannot
be cuspidal or septic-plus-cusp, proving the classification.

The complete source \texttt{classify\_octic\_33.m} exports the actual
Picard relation matrix, the rational cuspidal generators, all known
class sets and all $64660$ tested classes. The independent program
\texttt{replay\_class\_arithmetic.py} checks every group membership
and the exact equality of the remaining $105$ classes with the
$100+5$ known classes, using integer lattices.

The known octic elliptic curve has full torsion $C_{33}$:
good unramified odd-prime reductions have orders $5760579$ over
$\mathbb F_{7^8}$ and $815675619$ over $\mathbb F_{13^8}$,
whose gcd is $33$.
Finally, a point of $X_1(33)$ defined over $K$ yields a noncuspidal
rational trace divisor of degree $8$, allowing multiplicities if
its field of definition is smaller than $K$.
The classification forces this trace to be one of the known closed
octic divisors. Its field of definition is therefore all of $K$.
The fine moduli interpretation identifies the elliptic curve with
the known model over that field, so its full torsion is $C_{33}$.
\end{proof}
